\chapter{Ориентационный функтор Мориты для хопфовой линии}

Предыдущий гейт построил хопфову линию $L$ с $c_1(L)=1$ и показал, что
пара $E\otimes L$, $E^*\otimes L^*$ совместима с KO6. Оставался один
стоп-вопрос: является ли назначение $L/L^*$ направленным стрелкам
следствием родителя или очередным ручным знаком.

Настоящий гейт проверяет три возможных источника ориентации. Неединственная
трёхуровневая высота уже исключена предыдущим аудитом. Полярная
декомпозиция отдельного элемента $E$ также зависит от самого элемента и
не задаёт общую ориентацию пространства переходов. Остаётся структура
двух углов связывающей алгебры.

\section{Каноническая градуировка связывающей алгебры}

Связывающий родитель имеет вид
\begin{equation}
 \mathcal L(E)=
 \begin{pmatrix}
  M_{20}&E\\
  E^*&M_{15}
 \end{pmatrix}
 \simeq M_{35}.
\end{equation}
Обозначим угловые проекторы через $p_{20}$ и $p_{15}$. Они определяют
самосопряжённую инволюцию
\begin{equation}
 \Gamma_{\rm link}=p_{20}-p_{15},
 \qquad
 \Gamma_{\rm link}^2=I_{35}.
 \label{eq:v5-hopf-link-grading}
\end{equation}

Для однородного элемента $X$ определим степень равенством
\begin{equation}
 [\Gamma_{\rm link},X]=2\deg(X)X.
 \label{eq:v5-hopf-link-degree}
\end{equation}
Прямое блочное умножение даёт
\begin{equation}
 \deg(E)=+1,
 \qquad
 \deg(E^*)=-1,
 \qquad
 \deg(M_{20}\oplus M_{15})=0.
 \label{eq:v5-hopf-link-degrees}
\end{equation}

Это не новая высота аффинной цепи. Градуировка
\eqref{eq:v5-hopf-link-grading} существует уже у всякой двухугловой
связывающей алгебры и различает источник и цель перехода.

\section{Функториальность степени}

Если произведение двух однородных элементов ненулевое, то
\begin{equation}
 \deg(XY)=\deg(X)+\deg(Y).
 \label{eq:v5-hopf-degree-additivity}
\end{equation}
Сопряжение обращает направление:
\begin{equation}
 \deg(X^*)=-\deg(X).
 \label{eq:v5-hopf-degree-star}
\end{equation}

Машинный аудит проверяет эти утверждения на всех ненулевых произведениях
матричных единиц связывающего представителя $M_{3+2}$. Получено $125$
ненулевых композиций и ни одного дефекта аддитивности.

Теперь хопфово назначение задаётся формулой
\begin{equation}
 \mathfrak T(X)=X\otimes L^{\otimes\deg(X)},
 \qquad
 L^{-1}:=L^*.
 \label{eq:v5-hopf-orientation-functor}
\end{equation}
Аддитивность степени делает $\mathfrak T$ мультипликативным. В частности,
\begin{align}
 (E\otimes L)(E^*\otimes L^*)&\subset M_{20},\\
 (E^*\otimes L^*)(E\otimes L)&\subset M_{15},
\end{align}
поскольку $L\otimes L^*$ канонически сокращается в тривиальную линию.

\begin{theorem}[Ориентационный функтор двух углов]
Градуировка $\Gamma_{\rm link}$ канонически определяет аддитивную степень
на связывающей алгебре. Назначение
$E\mapsto E\otimes L$, $E^*\mapsto E^*\otimes L^*$ совместимо с
композицией и $*$-операцией и не требует отдельного знакового селектора.
\end{theorem}

\section{Почему углы нельзя незаметно переставить}

Для двух углов одинакового ранга оставалась бы внешняя симметрия их
перестановки. Здесь
\begin{equation}
 \rank p_{20}=20,
 \qquad
 \rank p_{15}=15.
\end{equation}
Унитарное сопряжение сохраняет ранг проектора, поэтому внутри $M_{35}$ не
существует внутреннего автоморфизма, меняющего $p_{20}$ и $p_{15}$
местами. Кроме того, их роли уже независимо зафиксированы как семейное и
наблюдаемое чтения.

Глобальная замена $\Gamma_{\rm link}\mapsto-\Gamma_{\rm link}$ не является
новой непрерывной свободой. Она одновременно меняет $E\leftrightarrow E^*$
и $L\leftrightarrow L^*$ и реализуется сопряжением/вещественной структурой.

\section{Выбор двух подъёмов проекторной оси}

Для точечного проекторного дефекта поле
\begin{equation}
 P=nn^T
\end{equation}
имеет два ориентированных подъёма $n$ и $-n$. На окружающей сфере $S^2$
они задают сопряжённые хопфовы линии
\begin{equation}
 c_1(L_n)=+1,
 \qquad
 c_1(L_{-n})=-1.
\end{equation}
Ориентационный функтор фиксирует
\begin{equation}
 \deg(E)=+1\longmapsto L_n,
 \qquad
 \deg(E^*)=-1\longmapsto L_{-n}=L_n^*.
 \label{eq:v5-hopf-arrow-lift-selection}
\end{equation}

Важна область применимости. Окружающая дефект сфера односвязна, поэтому
подъём карты в $\mathbb{RP}^2$ существует и имеет ровно две глобальные
ветви. Функтор выбирает их парно, не вводя локальных разрезов или
дополнительных наблюдаемых состояний.

\section{Совместимость с KO6 и следом}

Вещественная структура меняет направление перехода и комплексно
сопрягает коэффициентную линию:
\begin{equation}
 J(E\otimes L)J^{-1}=E^*\otimes L^*.
 \label{eq:v5-hopf-J-oriented-pair}
\end{equation}
Поэтому полный KO6-модуль сохраняется. Ориентированная физическая половина
видит $c_1=+1$, сопряжённая --- $c_1=-1$.

Нормированный след $M_{35}$ также не меняется. Линия имеет комплексный
ранг один, а во всех диагональных произведениях $L$ сокращается с $L^*$.
Новый относительный секторный коэффициент не возникает.

\section{Отношение к прежнему запрету высоты}

Предыдущий гейт высоты доказал, что KO6 и единичные перепады не выбирают
между трёхуровневой цепью и конфигурацией со стоком. Этот результат
остаётся в силе. Нынешняя ориентация возникает на другом уровне: не внутри
цепи $4\to3\to3$, а между двумя неравноранговыми углами
$M_{20}$ и $M_{15}$.

Следовательно, новая конструкция обходит прежнюю неоднозначность, но не
опровергает её.

\section{Вердикт и следующий предел}

Недостающий ориентационный функтор выведен из связывающего родителя. Это
закрывает топологический мост
\begin{equation}
 \boxed{
 E/E^*\longrightarrow L/L^*
 \longrightarrow c_1=\pm1}
\end{equation}
без нового $SU(2)_F$, без второй фиксированной двойки и без изменения
следа.

Но физическая частица ещё не построена. Хопфова связность пока является
канонической связностью коэффициентной линии на сфере вокруг дефекта. Нужно
вывести один пространственный суперсвязностный функционал, в котором её
кривизна одновременно:
\begin{enumerate}
 \item делает энергию проекторного ежа конечной;
 \item входит в скрученный оператор Дирака с индексом один;
 \item имеет относительную нормировку, фиксированную следом родителя.
\end{enumerate}

\begin{protocol}
\begin{enumerate}
 \item градуировка двух углов: \textbf{канонична};
 \item степень $E/E^*$: \textbf{$+1/-1$};
 \item аддитивность при композиции: \textbf{доказана};
 \item обращение степени при $*$: \textbf{доказано};
 \item назначение $L/L^*$: \textbf{функториально};
 \item углы $20/15$ внутренне переставимы: \textbf{нет};
 \item KO6: \textbf{сохранён};
 \item нормированный след: \textbf{сохранён};
 \item новый конечный модуль: \textbf{не введён};
 \item топологический ориентационный мост: \textbf{пройден};
 \item конечная энергия и единое действие: \textbf{открыты};
 \item физическое замыкание: \textbf{не достигнуто}.
\end{enumerate}
\end{protocol}

Машинный протокол сохранён в
\path{s2t_v5_hopf_line_morita_orientation_functor_gate_results.json}.